\documentclass{standalone}
\usepackage{tikzplot}

\begin{document}

\begin{tikzpicture}
  \tikzmath{
    % 5 点的齐次坐标
    \Pa{1,1} = 0; \Pa{2,1} = 0; \Pa{3,1} = 1; 
    \Pb{1,1} = 2; \Pb{2,1} = 0; \Pb{3,1} = 1; 
    \Pc{1,1} = 2; \Pc{2,1} = 2; \Pc{3,1} = 1; 
    \Pd{1,1} = 0; \Pd{2,1} = 4; \Pd{3,1} = 1; 
    \Pe{1,1} = -2; \Pe{2,1} = 2; \Pe{3,1} = 1; 
  }
  \tikzset{
    conics/define/five points={\Pa,\Pb,\Pc,\Pd,\Pe,\A},
    mat/stdout={\A,3,3},
  }

  \axes[grid] (-5:5,-5:5);
  \conic[draw=teal] (\A);

  \coordinate (P1) at (\Pa{1,1},\Pa{2,1});
  \coordinate (P2) at (\Pb{1,1},\Pb{2,1});
  \coordinate (P3) at (\Pc{1,1},\Pc{2,1});
  \coordinate (P4) at (\Pd{1,1},\Pd{2,1});
  \coordinate (P5) at (\Pe{1,1},\Pe{2,1});

  \foreach \p/\placement in {P1/above, P2/above, P3/above, P4/above, P5/above}{
    \fill[red] (\p) circle (1.5pt);
    \draw (\p) node[\placement] {$\p$};
  }
\end{tikzpicture}

\begin{tikzpicture}
  \tikzmath{
    % 5 点的齐次坐标
    \Pa{1,1} = -2; \Pa{2,1} = 0; \Pa{3,1} = 1; 
    \Pb{1,1} = -1; \Pb{2,1} = 1; \Pb{3,1} = 1; 
    \Pc{1,1} = -2; \Pc{2,1} = 4; \Pc{3,1} = 1; 
    \Pd{1,1} = 2; \Pd{2,1} = 0; \Pd{3,1} = 1; 
    \Pe{1,1} = 3; \Pe{2,1} = 1; \Pe{3,1} = 1; 
  }
  \tikzset{
    conics/define/five points={\Pa,\Pb,\Pc,\Pd,\Pe,\A},
    mat/stdout={\A,3,3},
  }

  \clip (-6,-6) rectangle (6,6);
  \axes[grid] (-5:5,-5:5);
  \conic[draw=teal] (\A);

  \coordinate (P1) at (\Pa{1,1},\Pa{2,1});
  \coordinate (P2) at (\Pb{1,1},\Pb{2,1});
  \coordinate (P3) at (\Pc{1,1},\Pc{2,1});
  \coordinate (P4) at (\Pd{1,1},\Pd{2,1});
  \coordinate (P5) at (\Pe{1,1},\Pe{2,1});

  \foreach \p/\placement in {P1/above, P2/above, P3/above, P4/above, P5/above}{
    \fill[red] (\p) circle (1.5pt);
    \draw (\p) node[\placement] {$\p$};
  }
\end{tikzpicture}

\begin{tikzpicture}
  \tikzmath{
    % 5 点的齐次坐标
    \Pa{1,1} = 1; \Pa{2,1} = 2; \Pa{3,1} = 1; 
    \Pb{1,1} = 4; \Pb{2,1} = 2; \Pb{3,1} = 1; 
    \Pc{1,1} = 1; \Pc{2,1} = -2; \Pc{3,1} = 1; 
    \Pd{1,1} = 4; \Pd{2,1} = 4; \Pd{3,1} = 1; 
    \Pe{1,1} = 5; \Pe{2,1} = 4; \Pe{3,1} = 1; 
  }
  \tikzset{
    conics/define/five points={\Pa,\Pb,\Pc,\Pd,\Pe,\A},
    mat/stdout={\A,3,3},
  }

  \axes[grid] (-3:6,-3:6);
  \conic[draw=teal,domain=-5:5] (\A);

  \coordinate (P1) at (\Pa{1,1},\Pa{2,1});
  \coordinate (P2) at (\Pb{1,1},\Pb{2,1});
  \coordinate (P3) at (\Pc{1,1},\Pc{2,1});
  \coordinate (P4) at (\Pd{1,1},\Pd{2,1});
  \coordinate (P5) at (\Pe{1,1},\Pe{2,1});

  \foreach \p/\placement in {P1/above, P2/above, P3/above, P4/above, P5/above}{
    \fill[red] (\p) circle (1.5pt);
    \draw (\p) node[\placement] {$\p$};
  } 
\end{tikzpicture}

\begin{tikzpicture}
  \tikzmath{
    % 5 点的齐次坐标
    \Pa{1,1} = 0; \Pa{2,1} = -2; \Pa{3,1} = 1; 
    \Pb{1,1} = 2; \Pb{2,1} = -2; \Pb{3,1} = 1; 
    \Pc{1,1} = 2; \Pc{2,1} = 0; \Pc{3,1} = 1; 
    \Pd{1,1} = 0; \Pd{2,1} = 2; \Pd{3,1} = 1; 
    \Pe{1,1} = -2; \Pe{2,1} = 0; \Pe{3,1} = 1; 
    \Pf{1,1} = -3; \Pf{2,1} = -3; \Pf{3,1} = 1;
  }
  \tikzset{
    conics/define/five points={\Pa,\Pb,\Pc,\Pd,\Pe,\A},
    mat/stdout={\A,3,3},
    conics/define/five points={\Pa,\Pb,\Pc,\Pd,\Pf,\B},
    mat/stdout={\B,3,3},
  }

  \clip (-6,-6) rectangle (6,6);
  \axes[grid] (-5:5,-5:5);
  \conic[draw=teal] (\A);
  \conic[draw=magenta,domain=-2.5:2.5] (\B);

  \coordinate (P1) at (\Pa{1,1},\Pa{2,1});
  \coordinate (P2) at (\Pb{1,1},\Pb{2,1});
  \coordinate (P3) at (\Pc{1,1},\Pc{2,1});
  \coordinate (P4) at (\Pd{1,1},\Pd{2,1});
  \coordinate (P5) at (\Pe{1,1},\Pe{2,1});
  \coordinate (P6) at (\Pf{1,1},\Pf{2,1});

  \foreach \p/\placement in {P1/below, P2/below right, P3/above right, P4/above, P5/below left, P6/below}{
    \fill[red] (\p) circle (1.5pt);
    \draw (\p) node[\placement] {$\p$};
  }
\end{tikzpicture}

\end{document}